% GENERATED FILE. Edit canonical lessons, then run npm run generate:books.
% canonical-source-hash: fnv1a64:3c91c2570acd600a
% canonical-lessons: KA-C215-mulaka, KA-C215-prakara, KA-C215-bagge, KA-C215-sampradaya, KA-C215-samskrti, KA-R215-first-pass-words-from-chapters-207-210, KA-R215-second-pass-words-from-chapters-211-215

\chapter{Through, According to, About, Tradition, Culture}
\label{ch:ka-through-according-to-about-tradition-culture}
\hlchaptermodality{215}

\begin{chapteropening}
\textbf{By the end of this chapter:} \emph{I can say through, according to, about, a tradition and culture.}

Complete the last lesson of chapter 215: I can say through, according to, about, a tradition and culture.
\end{chapteropening}

\section[\texorpdfstring{\kn{ಮೂಲಕ} (mūlaka)}{ಮೂಲಕ (mūlaka)}]{\kn{ಮೂಲಕ} (mūlaka) — through, by means of}

\label{lesson:KA-C215-mulaka}



\begin{quote}
Before the new one: say the Kannada for a cross street, then the Kannada for a fork.
\end{quote}

\subsection*{You'll want to know: \kn{ಮೂಲಕ}}
\textbf{\kn{ಮೂಲಕ}} — \emph{mūlaka} — \textquotedblleft{}through, by means of\textquotedblright{}.

\begin{grammarlens}[title={What you've built}]
One more word for this chapter.
\end{grammarlens}

\subsection*{Your turn}
\begin{itemize}
\raggedright
  \item \emph{Say it:} \emph{mūlaka}
  \item \emph{Say it:} \emph{mūlaka}, once more
  \item \emph{Say it:} say \emph{mūlaka}
  \item \emph{Recall:} say the Kannada for a cross street, then the Kannada for a fork, then say \emph{mūlaka} again
\end{itemize}

\begin{tcolorbox}[breakable,colback=teal!4,colframe=teal!35!black,title={Before you move on}]
What does \kn{ಮೂಲಕ} mean? (\textquotedblleft{}Through, by means of\textquotedblright{}.) Say it once more.
\end{tcolorbox}
\section[\texorpdfstring{\kn{ಪ್ರಕಾರ} (prakāra)}{ಪ್ರಕಾರ (prakāra)}]{\kn{ಪ್ರಕಾರ} (prakāra) — according to}

\label{lesson:KA-C215-prakara}



\begin{quote}
Before the new one: say the Kannada for a fork, then the Kannada for through.
\end{quote}

\subsection*{You'll want to know: \kn{ಪ್ರಕಾರ}}
\textbf{\kn{ಪ್ರಕಾರ}} — \emph{prakāra} — \textquotedblleft{}according to\textquotedblright{}.

\begin{grammarlens}[title={What you've built}]
One more word for this chapter.
\end{grammarlens}

\subsection*{Your turn}
\begin{itemize}
\raggedright
  \item \emph{Say it:} \emph{prakāra}
  \item \emph{Say it:} \emph{prakāra}, once more
  \item \emph{Say it:} say \emph{prakāra}
  \item \emph{Recall:} say the Kannada for a fork, then the Kannada for through, then say \emph{prakāra} again
\end{itemize}

\begin{tcolorbox}[breakable,colback=teal!4,colframe=teal!35!black,title={Before you move on}]
What does \kn{ಪ್ರಕಾರ} mean? (\textquotedblleft{}According to\textquotedblright{}.) Say it once more.
\end{tcolorbox}
\section[\texorpdfstring{\kn{ಬಗ್ಗೆ} (bagge)}{ಬಗ್ಗೆ (bagge)}]{\kn{ಬಗ್ಗೆ} (bagge) — about, concerning}

\label{lesson:KA-C215-bagge}



\begin{quote}
Before the new one: say the Kannada for through, then the Kannada for according to.
\end{quote}

\subsection*{You'll want to know: \kn{ಬಗ್ಗೆ}}
\textbf{\kn{ಬಗ್ಗೆ}} — \emph{bagge} — \textquotedblleft{}about, concerning\textquotedblright{}.

\begin{grammarlens}[title={What you've built}]
One more word for this chapter.
\end{grammarlens}

\subsection*{Your turn}
\begin{itemize}
\raggedright
  \item \emph{Say it:} \emph{bagge}
  \item \emph{Say it:} \emph{bagge}, once more
  \item \emph{Say it:} say \emph{bagge}
  \item \emph{Recall:} say the Kannada for through, then the Kannada for according to, then say \emph{bagge} again
\end{itemize}

\begin{tcolorbox}[breakable,colback=teal!4,colframe=teal!35!black,title={Before you move on}]
What does \kn{ಬಗ್ಗೆ} mean? (\textquotedblleft{}About, concerning\textquotedblright{}.) Say it once more.
\end{tcolorbox}
\section[\texorpdfstring{\kn{ಸಂಪ್ರದಾಯ} (sampradāya)}{ಸಂಪ್ರದಾಯ (sampradāya)}]{\kn{ಸಂಪ್ರದಾಯ} (sampradāya) — a tradition}

\label{lesson:KA-C215-sampradaya}



\begin{quote}
Before the new one: say the Kannada for according to, then the Kannada for about.
\end{quote}

\subsection*{You'll want to know: \kn{ಸಂಪ್ರದಾಯ}}
\textbf{\kn{ಸಂಪ್ರದಾಯ}} — \emph{sampradāya} — \textquotedblleft{}a tradition\textquotedblright{}.

\begin{grammarlens}[title={What you've built}]
One more word for this chapter.
\end{grammarlens}

\subsection*{Your turn}
\begin{itemize}
\raggedright
  \item \emph{Say it:} \emph{sampradāya}
  \item \emph{Say it:} \emph{sampradāya}, once more
  \item \emph{Say it:} say \emph{sampradāya}
  \item \emph{Recall:} say the Kannada for according to, then the Kannada for about, then say \emph{sampradāya} again
\end{itemize}

\begin{tcolorbox}[breakable,colback=teal!4,colframe=teal!35!black,title={Before you move on}]
What does \kn{ಸಂಪ್ರದಾಯ} mean? (\textquotedblleft{}A tradition\textquotedblright{}.) Say it once more.
\end{tcolorbox}
\section[\texorpdfstring{\kn{ಸಂಸ್ಕೃತಿ} (saṃskṛti)}{ಸಂಸ್ಕೃತಿ (saṃskṛti)}]{\kn{ಸಂಸ್ಕೃತಿ} (saṃskṛti) — culture}

\label{lesson:KA-C215-samskrti}



\begin{quote}
Before the new one: say the Kannada for about, then the Kannada for a tradition.
\end{quote}

\subsection*{You'll want to know: \kn{ಸಂಸ್ಕೃತಿ}}
\textbf{\kn{ಸಂಸ್ಕೃತಿ}} — \emph{saṃskṛti} — \textquotedblleft{}culture\textquotedblright{}.

\begin{grammarlens}[title={What you've built}]
One more word for this chapter.
\end{grammarlens}

\subsection*{Your turn}
\begin{itemize}
\raggedright
  \item \emph{Say it:} \emph{saṃskṛti}
  \item \emph{Say it:} \emph{saṃskṛti}, once more
  \item \emph{Say it:} say \emph{saṃskṛti}
  \item \emph{Recall:} say the Kannada for about, then the Kannada for a tradition, then say \emph{saṃskṛti} again
\end{itemize}

\begin{tcolorbox}[breakable,colback=teal!4,colframe=teal!35!black,title={Before you move on}]
What does \kn{ಸಂಸ್ಕೃತಿ} mean? (\textquotedblleft{}Culture\textquotedblright{}.) Say it once more.
\end{tcolorbox}
\section[(dialogue)]{Words from chapters 207-210 — a first pass}

\label{lesson:KA-R215-first-pass-words-from-chapters-207-210}



\begin{quote}
Say the words for a tradition and culture.
\end{quote}

\subsection*{Your turn}
Say these aloud:

\begin{itemize}
\raggedright
  \item to sniff, to hang up, to pound, to stack, to collect
  \item the body, an injection, an ointment, an illness, a patient
  \item the mind, disappointment, fun, yoga, meditation
  \item finally, usually, immediately, free time, Ugadi
\end{itemize}

\begin{tcolorbox}[breakable,colback=teal!4,colframe=teal!35!black,title={Before you move on}]
To sniff? (\textbf{\kn{ಮೂಸು}}, \emph{mūsu}.) To hang up? (\textbf{\kn{ನೇತುಹಾಕು}}, \emph{nētuhāku}.) To pound? (\textbf{\kn{ಕುಟ್ಟು}}, \emph{kuṭṭu}.) To stack? (\textbf{\kn{ಪೇರಿಸು}}, \emph{pērisu}.) To collect? (\textbf{\kn{ಸಂಗ್ರಹಿಸು}}, \emph{saṅgrahisu}.) The body? (\textbf{\kn{ದೇಹ}}, \emph{dēha}.) An injection? (\textbf{\kn{ಚುಚ್ಚುಮದ್ದು}}, \emph{cuccumaddu}.) An ointment? (\textbf{\kn{ಮುಲಾಮು}}, \emph{mulāmu}.) An illness? (\textbf{\kn{ಕಾಯಿಲೆ}}, \emph{kāyile}.) A patient? (\textbf{\kn{ರೋಗಿ}}, \emph{rōgi}.) The mind? (\textbf{\kn{ಮನಸ್ಸು}}, \emph{manassu}.) Disappointment? (\textbf{\kn{ನಿರಾಸೆ}}, \emph{nirāse}.) Fun? (\textbf{\kn{ತಮಾಷೆ}}, \emph{tamāṣe}.) Yoga? (\textbf{\kn{ಯೋಗ}}, \emph{yōga}.) Meditation? (\textbf{\kn{ಧ್ಯಾನ}}, \emph{dhyāna}.) Finally? (\textbf{\kn{ಕೊನೆಗೆ}}, \emph{konege}.) Usually? (\textbf{\kn{ಸಾಮಾನ್ಯವಾಗಿ}}, \emph{sāmānyavāgi}.) Immediately? (\textbf{\kn{ತಕ್ಷಣ}}, \emph{takṣaṇa}.) Free time? (\textbf{\kn{ಬಿಡುವು}}, \emph{biḍuvu}.) Ugadi? (\textbf{\kn{ಯುಗಾದಿ}}, \emph{yugādi}.)
\end{tcolorbox}
\section[(dialogue)]{Words from chapters 211-215 — a second pass}

\label{lesson:KA-R215-second-pass-words-from-chapters-211-215}



\begin{quote}
Say the words for regularly and so far.
\end{quote}

\subsection*{Your turn}
Say these aloud:

\begin{itemize}
\raggedright
  \item regularly, so far, until, for now, a period
  \item a port, a factory, a cinema hall, a ground, a direction
  \item a fort, a wedding hall, a gym, a stadium, a museum
  \item an edge, a tip, the bottom, a cross street, a fork
  \item through, according to, about, a tradition, culture
\end{itemize}

\begin{tcolorbox}[breakable,colback=teal!4,colframe=teal!35!black,title={Before you move on}]
Regularly? (\textbf{\kn{ನಿಯಮಿತವಾಗಿ}}, \emph{niyamitavāgi}.) So far? (\textbf{\kn{ಇಲ್ಲಿಯವರೆಗೆ}}, \emph{illiyavarege}.) Until? (\textbf{\kn{ವರೆಗೆ}}, \emph{varege}.) For now? (\textbf{\kn{ಸದ್ಯಕ್ಕೆ}}, \emph{sadyakke}.) A period? (\textbf{\kn{ಅವಧಿ}}, \emph{avadhi}.) A port? (\textbf{\kn{ಬಂದರು}}, \emph{bandaru}.) A factory? (\textbf{\kn{ಕಾರ್ಖಾನೆ}}, \emph{kārkhāne}.) A cinema hall? (\textbf{\kn{ಚಿತ್ರಮಂದಿರ}}, \emph{citramandira}.) A ground? (\textbf{\kn{ಮೈದಾನ}}, \emph{maidāna}.) A direction? (\textbf{\kn{ದಿಕ್ಕು}}, \emph{dikku}.) A fort? (\textbf{\kn{ಕೋಟೆ}}, \emph{kōṭe}.) A wedding hall? (\textbf{\kn{ಕಲ್ಯಾಣ} \kn{ಮಂಟಪ}}, \emph{kalyāṇa maṇṭapa}.) A gym? (\textbf{\kn{ಜಿಮ್}}, \emph{jim}.) A stadium? (\textbf{\kn{ಕ್ರೀಡಾಂಗಣ}}, \emph{krīḍāṅgaṇa}.) A museum? (\textbf{\kn{ವಸ್ತುಸಂಗ್ರಹಾಲಯ}}, \emph{vastusaṅgrahālaya}.) An edge? (\textbf{\kn{ಅಂಚು}}, \emph{añcu}.) A tip? (\textbf{\kn{ತುದಿ}}, \emph{tudi}.) The bottom? (\textbf{\kn{ತಳ}}, \emph{taḷa}.) A cross street? (\textbf{\kn{ಅಡ್ಡರಸ್ತೆ}}, \emph{aḍḍaraste}.) A fork? (\textbf{\kn{ಕವಲು}}, \emph{kavalu}.) Through? (\textbf{\kn{ಮೂಲಕ}}, \emph{mūlaka}.) According to? (\textbf{\kn{ಪ್ರಕಾರ}}, \emph{prakāra}.) About? (\textbf{\kn{ಬಗ್ಗೆ}}, \emph{bagge}.) A tradition? (\textbf{\kn{ಸಂಪ್ರದಾಯ}}, \emph{sampradāya}.) Culture? (\textbf{\kn{ಸಂಸ್ಕೃತಿ}}, \emph{saṃskṛti}.)
\end{tcolorbox}
